% 中文摘要与英文摘要分别成页
\clearpage
\phantomsection
\addcontentsline{toc}{section}{摘要}
\begin{center}
  {\heiti\zihao{2} 摘\quad 要}
\end{center}
\vspace{1.2em}

个性化辅导排课需要同时协调教师资质、班级可用时间、教室容量、设备条件
及临时变更等多类约束。依赖电子表格和人工核对的传统方式难以在业务变化
频繁时稳定保证无冲突，也不易解释不同排课方案之间的取舍。为此，本课题
拟研究并实现一套面向 K12 一对一及小班辅导场景的智能排课决策系统
ClassMind。

课题以约束规划和 OR-Tools CP-SAT 为核心，将“课次完整安排、教师/教室/
班级同一时段不冲突、教师资质匹配、教室容量及设备满足”等条件建模为硬
约束，将学生偏好、教师偏好、资源利用、日负载均衡和调课变更数建模为可
配置软目标。系统拟采用“业务数据层—约束求解层—独立校验层—应用交互层”
四层架构，形成候选方案生成、评分解释、异常诊断和最小影响调课的闭环，
并为第二阶段接入飞书多维表格、消息卡片和审批流程预留接口。

目前已完成可运行原型，包括数据模型、CP-SAT 求解器、本地 REST API、
四个功能页面和自动化测试。后续研究将围绕真实规模数据下的求解性能、
多目标权重敏感性、局部重排稳定性、结果可解释性和飞书场景集成展开。
预期成果包括一套可复现实验的智能排课原型、约束与评价指标体系、飞书
集成方案以及完整的技术文档。

\vspace{1em}
\noindent\textbf{关键词：}智能排课；约束规划；CP-SAT；多目标优化；
最小影响调课；可解释决策；飞书应用

\clearpage
\phantomsection
\addcontentsline{toc}{section}{Abstract}
\begin{center}
  {\bfseries\zihao{2} Abstract}
\end{center}
\vspace{1.2em}

Personalized tutoring timetabling must coordinate teacher qualifications,
class availability, room capacity, equipment requirements, and frequent
schedule changes. Spreadsheet-based workflows rely heavily on manual checks,
making it difficult to guarantee conflict-free schedules or explain the
trade-offs among alternatives. This project proposes ClassMind, an intelligent
timetabling decision system for one-to-one and small-group K12 tutoring.

The project adopts constraint programming with OR-Tools CP-SAT. Lesson
completion, teacher-room-class exclusivity, qualification matching, capacity,
and equipment requirements are modeled as hard constraints. Student and
teacher preferences, resource utilization, workload balance, and schedule
changes are modeled as configurable soft objectives. The proposed architecture
contains four layers: business data, constraint solving, independent
validation, and application interaction. It supports candidate generation,
score-based explanation, infeasibility diagnosis, and minimum-impact
rescheduling, while reserving interfaces for Feishu Base, message cards, and
approval workflows.

A runnable prototype has been completed, including the data model, CP-SAT
solver, local REST APIs, four functional pages, and automated tests. Subsequent
research will evaluate scalability, weight sensitivity, rescheduling stability,
interpretability, and Feishu integration. Expected deliverables include a
reproducible scheduling prototype, a constraint and evaluation framework, a
Feishu integration plan, and complete technical documentation.

\vspace{1em}
\noindent\textbf{Keywords:} Intelligent Timetabling; Constraint Programming;
CP-SAT; Multi-objective Optimization; Minimum-impact Rescheduling;
Explainable Decision; Feishu Application
\clearpage
